%=========================================================================
% main_plan.tex - ESQUELETO / MAPA DEL PAPER DE REVISTA (RT-OPP)
%
% Este NO es el paper. Es el documento de trabajo que fija el orden y el
% contenido de cada seccion de MAIN_FINAL.tex.
%
% COMO USARLO:
%   - Cada bloque tiene un ID (S1, S3.2, ...) que corresponde 1:1 con una
%     seccion o subseccion de MAIN_FINAL.tex. Usa esos IDs para pedir
%     cambios: "mueve S3.3 a S4", "parte S5.4 en dos", etc.
%   - Campos de cada bloque:
%       Rol       -> por que esta en el paper (funcion argumental)
%       Contenido -> que se dice, en orden
%       Elementos -> ecuaciones / figuras / tablas que viven ahi
%   - \pendiente{...} marca decisiones o contenido abierto.
%   - \nota{...}      marca observaciones de estructura.
%   - \upd{...} (rojo) comentario de la revision del hilo conductor y
%     \sug{...} (verde) la sugerencia asociada. Se borran cuando se decide.
%   - Figuras F1..F8, tablas T1..T5 y el Algoritmo A1 tienen su propio mapa
%     en la Seccion 3.
%=========================================================================
\documentclass[11pt,a4paper]{article}

\usepackage{lmodern}
\usepackage[T1]{fontenc}
\usepackage[utf8]{inputenc}
\usepackage[spanish,es-tabla,es-noquoting]{babel}
\usepackage[margin=2.2cm]{geometry}
\usepackage{amsmath,amssymb}
\usepackage{booktabs}
\usepackage{tabularx}
\usepackage{enumitem}
\usepackage{xcolor}
\usepackage{tikz}
\usetikzlibrary{arrows.meta,positioning}
\usepackage[colorlinks=true,linkcolor=blue!60!black,urlcolor=blue!60!black]{hyperref}

\definecolor{IDblue}{RGB}{20,80,150}
\definecolor{Warn}{RGB}{180,60,0}
\definecolor{Note}{RGB}{90,90,90}
\definecolor{Done}{RGB}{20,110,60}
\definecolor{Soft}{RGB}{240,244,250}

% --- Macros de formato del esqueleto ------------------------------------
\newcommand{\bloque}[2]{%
  \vspace{2mm}\noindent
  {\color{IDblue}\textbf{[#1]}}\ \textbf{#2}\par\vspace{0.5mm}}
\newcommand{\campo}[1]{\noindent\textit{#1:}\ }
\newcommand{\pendiente}[1]{{\color{Warn}\textbf{[PENDIENTE]} #1}}
\newcommand{\nota}[1]{{\color{Note}\textbf{[NOTA]} #1}}
\newcommand{\listo}{{\color{Done}\textbf{[LISTA]}}\ }
% Comentarios de la revision del hilo (rojo) y sugerencias (verde)
\newcommand{\upd}[1]{\par\smallskip\noindent{\color{red!75!black}\small
    \textbf{[COMENTARIO]} #1}\par}
\newcommand{\sug}[1]{\par\noindent{\color{green!45!black}\small
    \textbf{[SUGERENCIA]} #1}\par\smallskip}

% Simbolo de porcentaje seguro en texto y en modo matematico (babel-espanol).
\newcommand{\pct}{\text{\char`\%}}
\newcommand{\gv}{\boldsymbol{\gamma}}

\setlist[itemize]{leftmargin=6mm,itemsep=0.4mm,topsep=0.8mm}
\setlength{\tabcolsep}{4pt}
% Columna de ancho fijo con salto de linea (mapa de secciones).
\newcolumntype{P}[1]{>{\raggedright\arraybackslash}p{#1}}

\title{\textbf{Esqueleto del paper de revista}\\[2mm]
\large Real-Time Loss-Minimizing Multipolar Optimal Pulse Patterns\\
for Grid-Connected T-Type Converters With TDD Compliance}
\author{Documento de trabajo --- no es el manuscrito}
\date{\today}

\begin{document}
\maketitle

%=========================================================================
\section*{0. Ficha del documento}

\noindent\begin{tabularx}{\textwidth}{lX}
\toprule
Archivo destino    & \texttt{MAIN\_FINAL.tex} (un solo archivo) \\
Formato            & IEEEtran; hoy con la opción \texttt{conference}
                     \pendiente{pasar a \texttt{journal} cuando se defina la revista} \\
Revista objetivo   & \pendiente{definir (TPEL / TIE / JESTPE...) y su límite de páginas} \\
Idioma             & Inglés, estilo IEEE \\
Estado             & Off-line completo; RT-OPP descrito; mediciones de
                     pérdidas en el prototipo de 10~kW; faltan los resultados
                     del RT-OPP (convergencia, transitorios, red distorsionada,
                     costo) \\
Secciones          & 7 (S1--S7) \\
Figuras / tablas   & 8 figuras (F1--F8), 5 tablas (T1--T5), 1 algoritmo (A1);
                     propuestas F9 (montaje) y T6 (equipos) \\
Referencias        & 29, embebidas como \texttt{thebibliography}, orden de cita \\
\bottomrule
\end{tabularx}

%=========================================================================
\section*{1. Hilo argumental}

El paper sostiene una sola cadena: diseño off-line que minimiza pérdidas
$\rightarrow$ refinamiento en línea que mantiene el cumplimiento
$\rightarrow$ evidencia.

\vspace{3mm}
\begin{center}
\begin{tikzpicture}[
  b/.style={draw,rounded corners=2pt,fill=Soft,align=center,
            font=\footnotesize,minimum height=10mm,text width=36mm,inner sep=2pt},
  a/.style={-{Latex[length=2mm]},thick,draw=IDblue!70}
]
\node[b] (s1) at (0,0)       {\textbf{S1} Introducción\\T-Type SiC, OPP, tiempo real};
\node[b] (s2) at (4.6,0)     {\textbf{S2} Formulación\\armónica multipolar};
\node[b] (s3) at (9.2,0)     {\textbf{S3} Modelo de\\pérdidas + LCL};
\node[b] (s4) at (13.8,0)    {\textbf{S4} Problema off-line\\+ esquema de control};
\node[b] (s5) at (13.8,-2.6) {\textbf{S5} RT-OPP\\(SQP, una iteración GN)};
\node[b] (s6) at (9.2,-2.6)  {\textbf{S6} Resultados\\simulación + medición};
\node[b] (s7) at (4.6,-2.6)  {\textbf{S7} Conclusiones};
\draw[a] (s1)--(s2); \draw[a] (s2)--(s3); \draw[a] (s3)--(s4);
\draw[a] (s4)--(s5); \draw[a] (s5)--(s6); \draw[a] (s6)--(s7);
\end{tikzpicture}
\end{center}

\noindent\textbf{Las frases que el paper tiene que dejar claras:}
\begin{enumerate}[leftmargin=6mm]
\item El margen de distorsión que abre la modulación multipolar no se gasta
      en bajar más el TDD: se \emph{reinvierte} en bajar pérdidas hasta el
      límite de IEEE~519 (off-line, S2--S4).
\item Ese patrón deja de ser óptimo entre puntos de la tabla y con red
      distorsionada; el RT-OPP lo refina en línea con costo acotado por
      muestra y recupera el cumplimiento, partiendo de la tabla (S5).
\item La ganancia en pérdidas se ve en todo el rango de operación y se
      confirma con mediciones (S6).
\end{enumerate}

\upd{Frase 2: el RT minimiza una FO que combina distorsión y pérdidas
(decisión del autor), y los límites entran como pesos adaptativos, no como
restricciones duras. Eso explica que quede bajo el límite.}
\sug{Frase 2: ``Ese patrón deja de ser óptimo entre puntos de la tabla y con
red distorsionada; el RT-OPP lo refina en línea minimizando distorsión y
pérdidas, con los límites como pesos adaptativos, y recupera el cumplimiento
con costo acotado por muestra.''}

%=========================================================================
\section*{2. Mapa de secciones}

{\small
\noindent\begin{tabularx}{\textwidth}{@{}lP{3.7cm}XP{3.6cm}P{1.9cm}@{}}
\toprule
\textbf{ID} & \textbf{Sección} & \textbf{Rol en el argumento} & \textbf{Ecs.} & \textbf{Fig./Tab.}\\
\midrule
S1 & Introduction                      & Contexto, brechas y contribuciones & -- & F1\\
S2 & Multipolar Harmonic Formulation   & Convierte el OPP en un NLP continuo & fourier\_gamma, tdd, niveles & --\\
S3 & Loss Estimation                   & Pérdidas en función de $\gv$; filtro & p\_con, t\_i, pcond\_final, p\_sw, fres & F2, F3, T1, T2\\
S4 & OPP Implementation Scheme         & Problema off-line y lazo de control & opp\_propuesto & T3\\
S5 & Real-Time Implementation          & Refinamiento en línea (SQP) & rt\_lin \dots\ rt\_theta & A1, T4\\
S6 & Simulation and Experimental Results & Evidencia & -- & F4--F9, T5, T6\\
S7 & Conclusions                       & Cierre y trabajo futuro & -- & --\\
\bottomrule
\end{tabularx}}

\vspace{1mm}
\nota{La topología (T-Type + SiC) ya no es una sección propia: vive en el
párrafo 2 de S1. F1 se ancla después de S1, pero se cita por primera vez en
S3.1.}

%=========================================================================
\section*{3. Mapa de figuras, tablas y algoritmo}

\noindent\begin{tabularx}{\textwidth}{lllXl}
\toprule
\textbf{ID} & \textbf{Etiqueta} & \textbf{Vive en} & \textbf{Contenido / estado} & \textbf{Ancho}\\
\midrule
F1 & \texttt{fig:ttype}          & tras S1 & \listo Topología T-Type y estados P/O/N & 2 col.\\
F2 & \texttt{fig:patron\_opp}    & S3.1 & \listo Patrón multipolar $d=5$, pares $(\gamma_i,\gamma_{i+d})$ & 1 col.\\
F3 & \texttt{fig:control}        & S3.2 & \listo Esquema de control v4 (rediseño del autor) & 2 col.\\
F4 & \texttt{fig:patron\_pulsos} & S6.1 & Ángulos vs.\ $M$ y patrón \pendiente{verificar $d=19$; dependencia en $\varphi$} & 2 col.\\
F5 & \texttt{fig:fft\_comparacion} & S6.2 & TDD vs.\ potencia y espectro \pendiente{aún DPWM1 6450~Hz y $d=17$} & 2 col.\\
F6 & \texttt{fig:perdidas}       & S6.3 & \listo Mapas de pérdidas del modelo en $(I_d,I_q)$: OPP, PD-PWM, SV-GDPWM & 2 col.\\
F7 & \texttt{fig:comp\_gdpwm}    & S6.3 & \listo Modelo vs.\ medición, OPP vs.\ SV-GDPWM & 2 col.\\
F8 & \texttt{fig:comp\_pwm}      & S6.3 & \listo Modelo vs.\ medición, OPP vs.\ PD-PWM & 2 col.\\
F9 & \texttt{fig:setup}         & S6.0 & \pendiente{Foto del banco con los equipos numerados + diagrama de conexiones y puntos de medición} & 2 col.\\
\midrule
T1 & \texttt{tab:estados\_cond}  & S3.1 & \listo Estados de conducción por pierna & 1 col.\\
T2 & \texttt{tab:transiciones\_zvs} & S3.2 & \listo Transiciones y ZVS & 1 col.\\
T3 & \texttt{tab:parametros}     & S4.1 & \listo Parámetros del sistema ($d=19$, 950~Hz) & 1 col.\\
T4 & \texttt{tab:rt\_param}      & S5 & \listo Parámetros del RT-OPP & 1 col.\\
T5 & \texttt{tab:comparacion}    & S6.3 & \listo Pérdidas totales en el punto nominal, modelo y medición & 1 col.\\
T6 & \texttt{tab:equipos}        & S6.0 & \pendiente{Equipos: fuente DC, red/emulador, OPAL-RT, analizador de potencia, sondas, osciloscopio} & 1 col.\\
\midrule
A1 & \texttt{alg:rtopp}          & S5 & \listo Refinamiento RT del OPP & 1 col.\\
\bottomrule
\end{tabularx}

\vspace{1mm}
\noindent\pendiente{Figuras citadas como \textbf{[X]} que aún no existen:
barrido de $\lambda_\ell$ (S5.4), recuperación del cumplimiento con red
distorsionada (S6.2) y convergencia del RT-OPP (S6.4 propuesta).}

%=========================================================================
\section*{4. Detalle por bloque}

\bloque{Abstract}{}
\campo{Contenido} Problema (baja $f_{sw}$ degrada la corriente) $\rightarrow$
OPP multipolar para T-Type SiC que minimiza pérdidas sujeto a IEEE~519
$\rightarrow$ modelos de pérdidas en ángulos virtuales y restricción de nivel
lineal $\rightarrow$ reducciones medias de 36/34\pct{} (modelo) y 33/8\pct{}
(medido, hasta 14\pct{}) frente a SV-GDPWM y PD-PWM.\\
\pendiente{El texto principal aún no menciona el RT-OPP; hay una \texttt{\textbackslash sug}
completa. El autor lo revisa al final.}

\bloque{S1}{Introduction}
\campo{Rol} Instalar el problema, ubicar las dos brechas y declarar el aporte.\\
\campo{Contenido}
\begin{itemize}
\item P1: alta potencia $\Rightarrow$ pérdidas relevantes; dos frentes:
      topología/dispositivos y modulación.
\item P2: T-Type + SiC frente al NPC (cruce 10--15~kHz con Si, se desplaza con
      SiC); transiciones de un nivel; objeción de bloqueo a 750~V.
\item P3: familias de modulación $\rightarrow$ SHE/SHM/SHC $\rightarrow$ OPP
      por programación no lineal; LCL dentro del problema.
\item P4: OPP multipolar; OPP con pérdidas o temperatura y ángulo virtual
      (todos NPC con IGCT) $\Rightarrow$ \textbf{brecha 1}: no hay modelo de
      pérdidas del T-Type en función de los ángulos, menos con SiC.
\item P5: OPP off-line pierden optimalidad entre puntos y con red
      distorsionada (IEEE~519); tres familias en línea; Tsekouras
      $\Rightarrow$ \textbf{brecha 2}: ningún OPP multipolar con pérdidas se
      refina en línea.
\item Contribuciones (4): modelo de pérdidas; formulación con restricción
      lineal de nivel; RT-OPP con una iteración GN; validación experimental.
\end{itemize}
\campo{Elementos} F1 (anclada aquí).

\upd{(a) Las contribuciones van en orden inverso al cuerpo (pérdidas antes
que formulación). (b) La contribución 2 no tiene brecha en S1: que
\texttt{koukoula\_junction} mantiene una restricción entera de nivel aparece
recién en S2.3. (c) Contribución 3: ``Hessian positive definite by
construction'' solo vale si $\lambda_u>2\lambda_\ell c$.}
\sug{(a) Orden: formulación, pérdidas, RT, validación experimental.
(b) En P4, después de citar el ángulo virtual: ``\dots\ although the level
constraint remains an integer one that must be reconstructed at every solver
iteration \texttt{\textbackslash cite\{koukoula\_junction\}}.'' (c) ``whose Hessian is positive
definite for a proximal weight above a bound given by the loss weight''.}

\bloque{S2*}{(nueva, propuesta) System Description}
\upd{La revisión propone juntar en una sección propia todo lo que define el
sistema y que hoy está repartido: topología (párrafo 2 de S1 y F1), estados
de conducción (T1, en S3.1), filtro LCL (S3.3, hoy dentro de ``Loss
Estimation''), parámetros (T3, en S4.1), y el esquema de control F3 con
el observador (hoy F3 flota en S3.2 y se cita desde S3.1 para definir
$\varphi$).}
\sug{Contenido: T-Type SiC y estados P/O/N (F1, T1); filtro LCL y
resonancia (\texttt{eq:fres}); parámetros (T3); cadena de control resumida:
$dq$ $\rightarrow$ $M,\delta$; $\varphi$ desde $P^\star,Q^\star$; observador de
armónicos; LUT $\rightarrow$ RT-OPP $\rightarrow$ generador (F3). Cierre:
``These elements define the quantities ($g_n$, $\varphi$, $I$) on which the
pattern design depends.'' La motivación de T-Type + SiC sigue en S1.}

\bloque{S2}{Multipolar Harmonic Formulation}
\campo{Rol} Convertir el OPP en un problema continuo y sin enteros.

\bloque{S2.1}{Multipolar Modulation and Choice of Symmetry}
\campo{Contenido} HWS frente a QWS (QWS bloquea las soluciones multipolares)
y frente a onda completa (componente continua).

\bloque{S2.2}{Virtual Angle and Fourier Coefficients}
\campo{Contenido} Ángulo virtual $\gamma_i$; coeficientes de Fourier sin
enteros; TDD normalizada a $I_{nom}$ con la ganancia $g_n$ del LCL; conjunto
$\mathcal{N}=\{1,5,7,\dots,49\}$.\\
\campo{Elementos} \texttt{eq:fourier\_gamma}, \texttt{eq:tdd}.

\bloque{S2.3}{Linear Level Constraints}
\campo{Contenido} Las dos desigualdades lineales que reemplazan
$-1\le u_i\le 1$; separación mínima $\epsilon_{sep}$ en tiempo real.\\
\campo{Elementos} \texttt{eq:niveles}.

\upd{S2.3: $\epsilon_{sep}=0.2^\circ$ se define como algo del RT, pero no
aparece en S5.5 ni en T4 (el RT usa barrera y $\tau_B$).}
\sug{Mover la frase de $\epsilon_{sep}$ a S5.5 (párrafo de ángulos) y
agregarla a T4, o borrarla.}

\bloque{S3}{Loss Estimation}
\campo{Rol} Escribir las pérdidas como funciones cerradas y diferenciables
de $\gv$.

\bloque{S3.1}{Conduction Losses}
\campo{Contenido} Estados P/O/N; potencia instantánea resistiva con
$\varphi$; pares $(\gamma_i,\gamma_{i+d})$ como pulsos; descomposición
($S_2+S_3$ todo el periodo, luego corrección por pulso); primitiva exacta $G$.\\
\campo{Elementos} \texttt{eq:p\_con}, \texttt{eq:t\_i},
\texttt{eq:pcond\_final}, T1, F2.

\upd{(a) $\varphi$ se define sin llamarlo ``load angle'' y apoyándose en F3,
que aún no se presenta. (b) La corriente aparece como $I$ (S3.1),
$I_{cond}$ ``corresponding to the modulation index'' (S3.2) e $I_{nom}$
(S5); en conexión a red la corriente la fijan $P,Q$, no $M$.}
\sug{(a) ``where $\varphi$, the load angle, is \dots''; si se crea S2*, F3 ya
estará presentada. (b) Un solo símbolo de amplitud de corriente, definido en
S2* o T3, y una frase en S4.1 que diga qué corriente usa cada punto
$(M,\varphi)$ de la tabla.}

\bloque{S3.2}{Switching Losses}
\campo{Contenido} Tensión conmutada $V_{dc}/2$; ZVS por dirección de la
corriente; un encendido por $\gamma_i$ y un apagado por $\gamma_{i+d}$.\\
\campo{Elementos} \texttt{eq:p\_sw}, T2; F3 anclada aquí.\\
\nota{F3 (esquema de control) flota en S3.2 pero pertenece a S4.2.}

\upd{F3 (esquema de control) está anclada aquí pero pertenece a S4.2 o a
S2*.}
\sug{Mover el bloque \texttt{figure*} de F3 a S2* (o al inicio de S4.2).}

\bloque{S3.3}{LCL Filter}
\campo{Contenido} Resonancia en el armónico 21; 19 y 23 los más
restrictivos; $L_1=L_2$ y $C$ mínimo.\\
\campo{Elementos} \texttt{eq:fres}.

\upd{El LCL no es una pérdida: lo usan $g_n$ (S2.2, con referencia hacia
adelante), el cálculo de $\varphi$ y $\mathbf{Q}$ (S5.2). También cita
\texttt{Section\textasciitilde\textbackslash ref\{sec:rt\}} hacia adelante.}
\sug{Mover S3.3 completo a S2*.}

\bloque{S4}{OPP Implementation Scheme}
\campo{Rol} Juntar los modelos en el problema off-line y mostrar cómo se usa
la solución.

\bloque{S4.1}{Optimization Problem}
\campo{Contenido} Problema por punto $(M,\varphi)$: mínimas pérdidas sujeto a
fundamental, niveles, orden, TDD, límites individuales y continuidad en $M$.
Contraste: off-line minimiza pérdidas con la distorsión como restricción; el
RT minimiza distorsión con las pérdidas como término secundario.\\
\campo{Elementos} \texttt{eq:opp\_propuesto}, T3.\\
\pendiente{Cómo se impuso la continuidad en $\varphi$ en Julia.}

\upd{(a) $\delta$ tiene dos significados: cota de continuidad en
\texttt{eq:opp\_propuesto} (sin definir en el texto) y fase del bloque
polar en S4.2. (b) El procedimiento de solución (1000 semillas, Ipopt, LUT
$9\times25$) está en el párrafo inicial de S6, pero es método. (c) Punto 2
del revisor: falta decir que aquí TDD y límites individuales son
restricciones duras y en el RT van en la FO (hay \texttt{\textbackslash sug}
en el paper).}
\sug{(a) Renombrar la cota a $\Delta\gamma_{\max}$ y definirla en el párrafo
que sigue a la ecuación. (b) Traer a S4.1 el procedimiento multi-start y la
construcción de la LUT $(M,\varphi)$; título posible ``Off-Line Pattern
Design''. (c) Aplicar la \texttt{\textbackslash sug} de S4.1 del paper.}

\bloque{S4.2}{Control and Pattern Generation}
\campo{Contenido} LUT $\gv(M,\varphi)$; lazo $dq$ $\rightarrow$ $M$ y
$\delta$; $\varphi$ desde $P^\star,Q^\star$ sin mediciones; observador de
armónicos; RT-OPP; generador desplazado por $\delta$; continuidad y
re-siembra con histéresis.\\
\campo{Elementos} F3.

\upd{(a) S4.2 adelanta el RT (``proximal term'', re-siembra con
histéresis) antes de S5; la re-siembra se explica 3 veces (S4.2, A1, S5.5) y
la motivación del RT 4 veces (S1, S4, S4.2, S5). (b) La continuidad se
argumenta solo en $M$; al re-sembrar en $\varphi$ los ángulos saltan y no se
dice cómo se acota el transitorio. (c) ``The LUT, now addressed by both $M$
and $\varphi$'' (S5) es resto de la versión de conferencia. (d) No se explica
por qué $M$ tiene rango estrecho (0.835--0.9) y $\varphi$ completo.
(e) Falta el observador (decisión del autor: agregarlo; hay
\texttt{\textbackslash sug} en el paper).}
\sug{(a) Dejar en S4.2 solo LUT y cadena de control; la re-siembra solo en
S5.5. (b) Una frase en S5.5: cómo se evita o acota el transitorio al cambiar
de celda en $\varphi$. (c) Borrar ``now''. (d) ``In a grid-connected
converter, $M$ is set mostly by the grid voltage, whereas $\varphi$ spans the
whole circle with $P^\star$ and $Q^\star$; the table is therefore dense in
$\varphi$ and narrow in $M$.''}

\bloque{S5}{Real-Time Implementation}
\campo{Rol} Refinar el patrón en línea con costo acotado por muestra.\\
\campo{Contenido} Por qué refinar (entre puntos de la tabla, red
distorsionada); SQP con una iteración GN por actualización y warm start;
restricciones como pesos adaptativos; la LUT da semilla y re-siembras.\\
\campo{Elementos} A1, T4.

\upd{La plataforma (OPAL-RT, $T_s=200~\mu$s) está en S6, pero es parte de
la implementación. ``OF'' se usa antes de definirse.}
\sug{Mover la frase de la plataforma y $T_s$ al párrafo inicial de S5;
escribir ``objective function (OF)'' en su primer uso.}

\bloque{S5.1}{Linearized Harmonic Model}
\campo{Contenido} $\mathbf{x}$ con los armónicos de $\mathcal{N}$; Jacobiano
exacto; el factor $1/n$ se cancela.\\
\campo{Elementos} \texttt{eq:rt\_lin}, \texttt{eq:rt\_jac}.

\bloque{S5.2}{TDD Term}
\campo{Contenido} Espectro de referencia $\mathbf{x}^\star$ (fundamental +
fondo de red estimado por el observador); TDD con fondo; matriz $\mathbf{Q}$
con pesos adaptativos; término de rizado 53--97.\\
\campo{Elementos} \texttt{eq:rt\_tdd}, \texttt{eq:rt\_Q}.

\bloque{S5.3}{Loss Terms}
\campo{Contenido} Conducción con signo ($t_i$) y conmutación en magnitud
($e_i$); sus linealizaciones.\\
\campo{Elementos} \texttt{eq:rt\_T}, \texttt{eq:rt\_pcond}, \texttt{eq:rt\_BP},
\texttt{eq:rt\_S}, \texttt{eq:rt\_psw}, \texttt{eq:rt\_BS}.

\upd{Dice que los modelos de S3 se ``reusan'', pero el de conmutación cambia
($2f_1$, $P_{sw,0}$, corrección por $R_G$).}
\sug{``\dots\ are reused with three adjustments: \dots''; una frase por
ajuste.}

\bloque{S5.4}{Local Quadratic Problem}
\campo{Contenido} Función objetivo cuadrática; ecuaciones normales;
$\mathbf{W}$ definida positiva si $\lambda_u>2\lambda_\ell c$; los residuos
cuadráticos subponderan la conmutación (las pérdidas actúan como
regularizador, $\lambda_\ell$ fija el nivel de TDD); costo computacional.\\
\campo{Elementos} \texttt{eq:rt\_fo}, \texttt{eq:rt\_W}.\\
\pendiente{Figura del barrido de $\lambda_\ell$ (Fig.~[X]); cifras de
\texttt{RTOPP\_D19}, regenerar con \texttt{RTOPP\_D19\_simple\_rip}.}

\upd{(a) Decisión del autor: no decir que las pérdidas solo actúan como
regularizador (hay \texttt{\textbackslash sug} en el paper). (b) El barrido de
$\lambda_\ell$ (Fig.~[X]) es un resultado. (c) Cita F6 (resultados) hacia
adelante. (d) Tres cifras para el cambio de pérdidas en RT: $<0.3$\pct{}
(S5.4), 0.4\pct{} (S6.2), 1\pct{} (\texttt{\textbackslash sug} de S6.4 y
S7).}
\sug{(b) Mover el barrido de $\lambda_\ell$ a la subsección de resultados RT
y dejar aquí solo el criterio para elegir $\lambda_\ell$. (c) Quitar la cita
a F6 o reemplazarla por la tabla de parámetros. (d) Unificar en una sola
cifra.}

\bloque{S5.5}{Constraint Handling}
\campo{Contenido} Límites IEEE~519 y TDD como pesos con integrador
logarítmico (lectura KKT); fundamental como igualdad + trim de Newton;
barrera logarítmica sobre el poliedro de ángulos; re-siembras.\\
\campo{Elementos} \texttt{eq:rt\_theta}.

\upd{Falta el transitorio al re-sembrar en $\varphi$ (ver S4.2).}
\sug{Una frase al final del párrafo de re-siembras.}

\bloque{S6}{Simulation and Experimental Results}
\campo{Rol} Mostrar que el diseño cumple y que la ganancia es real.\\
\campo{Contenido (párrafo inicial)} Solución en Julia/Ipopt; verificación en
PLECS; tabla RT $9\times25$; OPAL-RT OP4510 a 200~$\mu$s.\\
\pendiente{Describe aún la LUT de conferencia (1000 puntos, $d=17$, 8~min);
tiempo de ejecución [XX]; frase de PLECS frente a figuras ``sinplecs''.}

\upd{(a) El párrafo inicial mezcla método (multi-start, LUT, OPAL) con
resultados. (b) No se dice qué resultado es simulación, cuál HIL y cuál
medición; falta el montaje experimental. (c) El rango de operación aparece
de tres formas: potencia (F5), plano $(I_d,I_q)$ (F6--F8) y $M$ (F4).}
\sug{(a) Mover el método a S4.1/S5. (b) Ver el bloque nuevo S6.0. (c) Una frase que
relacione $|I|$, $P$ y $M$ (p.ej.\ el punto nominal $I_d=20$~A $\Leftrightarrow$
9.8~kW).}

\bloque{S6.0}{(nueva, propuesta) Experimental Setup}
\campo{Rol} Decir con qué se obtuvo cada resultado y por qué la medición de
pérdidas es confiable.\\
\campo{Contenido}
\begin{itemize}
\item Prototipo T-Type SiC de 10~kW, 750~V (módulo CHB011M12GM4T) con el
      filtro LCL de T3; conexión a red \textbf{[directa, transformador o
      emulador de red]}.
\item Alimentación DC \textbf{[fuente, rango]} y lado AC \textbf{[red,
      emulador o carga]}.
\item Control y modulación en OPAL-RT OP4510: RT-OPP a $T_s=200~\mu$s y
      generación de pulsos \textbf{[FPGA, resolución]}.
\item Medición de pérdidas: \textbf{[analizador de potencia, precisión]},
      método \textbf{[$P_{in}-P_{out}$ o calorimétrico]}, sondas de tensión y
      corriente, osciloscopio.
\item Qué resultado viene de cada plataforma: simulación (Julia, PLECS),
      HIL (OPAL-RT) y prototipo.
\end{itemize}
\campo{Elementos} F9 (foto + diagrama de conexiones), T6 (equipos).
\upd{Hoy el paper no describe el montaje. De esto dependen dos preguntas del
revisor: qué incluye la pérdida medida y si la precisión alcanza para las
diferencias de 0.4--0.5~W a carga ligera.}
\sug{F9 en dos paneles: (a) foto del banco con los equipos numerados;
(b) diagrama de conexiones con los puntos donde se miden $P_{in}$ y
$P_{out}$. T6 con columnas Equipo | Modelo | Uso | Precisión. Párrafo:
``The experimental setup is shown in Fig.~[F9] and the equipment is listed
in Table~[T6]. The 10~kW SiC T-Type prototype is supplied by [DC source] and
connected to [grid/emulator] through the LCL filter of Table~III. The
modulation and the RT-OPP run on an OPAL-RT OP4510 with $T_s=200~\mu$s. The
losses are obtained as [the difference between the DC input power and the AC
output power] measured by [power analyzer], whose accuracy of [XX] amounts to
[XX]~W at the nominal point.''}

\bloque{S6.1}{Optimized Angles and Pulse Pattern}
\campo{Contenido} Patrón multipolar en $M=0.8660$; conmutaciones cerca de
los cruces por cero; trayectorias continuas en $M$.\\
\campo{Elementos} F4.

\bloque{S6.2}{Harmonic Compliance}
\campo{Contenido} TDD en el límite en todo el rango; espectro dentro de los
límites individuales; RT-OPP: 1.7\pct{} en un punto de la tabla y
recuperación del cumplimiento en 13~ms con fondo armónico.\\
\campo{Elementos} F5.\\
\pendiente{F5 y su párrafo comparan con DPWM1 y $d=17$; cifra del 9\pct{}
sin ensayo definido.}

\upd{(a) Punto 2 del revisor: falta explicar por qué el RT queda bajo el
límite (hay \texttt{\textbackslash sug} en el paper). (b) $M=0.87$ no cae en la
grilla de 9 índices entre 0.835 y 0.9 (paso $\approx0.0081$), pero se le
llama ``table point''.}
\sug{(a) Aplicar la \texttt{\textbackslash sug} de S6.2 del paper. (b) Verificar
la grilla y usar un nodo real (p.ej.\ $M=0.8756$) o no llamarlo ``table
point''. Separar en dos subsecciones: cumplimiento off-line y resultados RT.}

\bloque{S6.3}{Loss Comparison}
\campo{Contenido} Benchmarks PD-PWM 2850~Hz y SV-GDPWM 4050~Hz (mínimo $m_f$
que cumple IEEE~519); mapas del modelo en $(I_d,I_q)$; modelo vs.\
medición en 8 puntos; tres efectos que explican la brecha frente a PD-PWM
(offset no modelado, carga ligera, rizado); mejora hasta 14.4/37.7\pct{}.\\
\campo{Elementos} F6, F7, F8, T5.\\
\pendiente{Montaje y método de medición; efectos no modelados [XX]; offset
mayor de SV-GDPWM; confirmar que se midió con el RT-OPP activo.}

\upd{(a) Las medias 6.8/10.1/10.3~W dan reducciones de 33/34\pct{}; el texto
dice 34/36\pct{}, que es la media de las reducciones punto a punto. (b) El
rizado aparece suelto en tres lugares: restricción off-line comentada,
término $\lambda_r$ del RT y brecha medida. (c) Pregunta probable: ¿es justo
comparar 950~Hz por dispositivo con portadoras de 2850 y 4050~Hz?}
\sug{(a) ``mean of the point-wise reductions''. (b) Una frase que conecte:
el off-line no restringe el rizado, el RT lo penaliza con $\lambda_r$, y el
rizado residual explica parte de la brecha medida. (c) Una frase con la
frecuencia de conmutación promedio por dispositivo de cada benchmark. Llevar
S6.3 al final de S6, después de los resultados RT.}

\bloque{S6.4}{(propuesta) Convergence of the Real-Time Optimizer}
\campo{Contenido} Hoy es una \texttt{\textbackslash sug}: convergencia desde
la semilla, punto entre celdas, red distorsionada, semilla perturbada.\\
\pendiente{Faltan figuras y datos; incluye transitorios ante escalones de
$M$ y $\varphi$ y costo por muestra.}

\upd{Debe juntar todos los resultados RT.}
\sug{Contenido: barrido de $\lambda_\ell$ (desde S5.4), cumplimiento RT y red
distorsionada (desde S6.2), convergencia, transitorios ante escalones de $M$ y
$\varphi$, costo por muestra.}

\bloque{S7}{Conclusions}
\campo{Contenido} Formulación continua y modelos de pérdidas; cifras de la
comparación (modelo y medición); el margen de distorsión convertido en
pérdidas.\\
\pendiente{Párrafo del RT-OPP (hay \texttt{\textbackslash sug}); frase de
trabajo futuro.}

\upd{(a) ``The optimizer drives the current distortion to the regulatory
limit'' choca con el RT (punto 2). (b) ``the 1000-point LUT \dots\ under
8~min'' es de la versión de conferencia. (c) Decisión del autor: el RT sí se
validó; ya no va como trabajo futuro (notas quitadas del paper).}
\sug{(a) ``The off-line optimizer drives \dots''. (b) Actualizar con la LUT
$(M,\varphi)$ y $d=19$. Orden del cierre: off-line $\rightarrow$ RT
$\rightarrow$ medición, igual que el cuerpo.}

%=========================================================================
\section*{5. Decisiones abiertas}

\begin{itemize}
\item \pendiente{Revista objetivo, límite de páginas y clase
      \texttt{journal}; biografías de autores con foto.}
\item \pendiente{Resultados del RT-OPP (S6.4) y figuras [X] de S5.4 y S6.2.}
\item \pendiente{F5 con los benchmarks y $d$ actuales.}
\item \pendiente{Párrafo inicial de S6 con la LUT $(M,\varphi)$ y $d=19$.}
\item \pendiente{Bug de signo en \texttt{P\_loss} del código (afecta solo
      el valor reportado).}
\item \pendiente{Referencias candidatas (\texttt{docs/REFERENCIAS.md}).}
\end{itemize}

%=========================================================================
\section*{6. Retroalimentación del hilo conductor}

\noindent Resumen de la revisión (detalle con líneas en
\texttt{docs/REVISION\_HILO.md}). Nada de esto está aplicado aún al paper.

\subsection*{6.1 Decisiones del autor sobre los puntos de severidad alta}
\begin{enumerate}[leftmargin=6mm]
\item \textbf{RT y pérdidas:} el RT sí minimiza pérdidas en línea. Se
      mantienen título y contribuciones; se cambia la frase de S5.4 que dice
      que las pérdidas solo actúan como regularizador.
\item \textbf{TDD 1.7\pct{} en RT frente a 5\pct{} off-line:} explicar que
      el off-line impone TDD y límites como restricciones duras y el RT los
      lleva en la FO, así que con el peso elegido no queda en el límite
      (S4.1, S6.2, S7).
\item \textbf{Validación del RT:} sí se validó; se quita como trabajo futuro
      de las conclusiones.
\item \textbf{Observador:} se agrega su descripción (S4.2 o S2*).
\end{enumerate}
\noindent Los demás hallazgos están como comentarios rojos con su
sugerencia verde en cada bloque de la Sección 4.

\subsection*{6.3 Esqueleto propuesto por la revisión}
\noindent\begin{tabularx}{\textwidth}{lX}
\toprule
\textbf{Nuevo} & \textbf{Contenido (IDs actuales)}\\
\midrule
S1 & Introduction; contribuciones en el orden del cuerpo \\
S2 & \emph{System Description} (nueva): topología + F1 + T1, LCL (S3.3),
     T3, F3 con la cadena de control y el observador \\
S3 & Formulación (S2 actual), sin $\epsilon_{sep}$ \\
S4 & Pérdidas (S3.1--S3.2 actual) \\
S5 & \emph{Off-Line Pattern Design}: S4.1 + multi-start + LUT
     $(M,\varphi)$ $9\times25$ + continuidad en $M$ y $\varphi$ \\
S6 & RT-OPP (S5 actual), con la cadena LUT$\rightarrow$semilla, observador,
     plataforma y $T_s$; re-siembra solo en S5.5; sin barrido de
     $\lambda_\ell$ \\
S7 & Results: 7.1 montaje; 7.2 ángulos (S6.1); 7.3 cumplimiento off-line;
     7.4 resultados RT (cumplimiento RT, barrido de $\lambda_\ell$,
     convergencia, transitorios, costo); 7.5 pérdidas (S6.3) \\
S8 & Conclusions: off-line $\rightarrow$ RT $\rightarrow$ medición \\
\bottomrule
\end{tabularx}

\subsection*{6.4 Preguntas probables del revisor}
\begin{enumerate}[leftmargin=6mm]
\item Si el RT baja la TDD a 1.7\pct{} sin cambiar pérdidas, ¿qué restricción
      está activa off-line?
\item ¿Se midió con el RT activo? ¿Dónde está la validación experimental del
      RT?
\item ¿Cómo funciona el observador y qué parte de los 13~ms le corresponde?
\item ¿Cómo se evitan los transitorios al re-sembrar en $\varphi$?
\item ¿Es justo comparar 950~Hz por dispositivo con portadoras de 2850 y
      4050~Hz?
\end{enumerate}

%=========================================================================
\section*{7. Espacio para tus comentarios}

\vspace{50mm}

\end{document}
